\documentclass{article}
\usepackage[margin=1in]{geometry}
\usepackage[skip=1em, indent=12pt]{parskip}
\usepackage{amsmath}
\usepackage{amssymb}
\usepackage{graphicx}
\usepackage{microtype}
\usepackage{textcomp, gensymb}
\usepackage{enumitem}
\begin{document}
\setcounter{section}{3}
\section{Chapter 4: Continuous-Time System Analysis using the Laplace Transform} 
\setcounter{subsection}{0}
\subsection{The Laplace Transform}
As a review, the bilateral and unilateral Laplace transforms may be expressed as 
\begin{align*}
	& X(s) = \int_{-\infty}^{\infty}x(t) e^{-st} dt \\
	& X(s) = \int_{0^-}^{\infty}x(t) e^{-st} dt,
\end{align*}
and the inverse Laplace transform may be expressed as 
\[
	x(t) = \frac{1}{2\pi j}\int_{c-j\infty}^{c+j\infty}X(s)e^{st}dt.
\]

To understand the Laplace transform, recall that \(s= \sigma + j\omega\), the complex frequency of the signal \(e^{st}\). Laplace transforms map a time-domain signal into the s-domain. In Circuits 2, we saw how transforming complex, second- or third-order circuits into the s-domain made it possible to solve them algebraically using KCL and KVL, rather than through a complex integro-differential equation representing its elements. Then, transforming them back into the time domain via the inverse Laplace transform, we could derive the solution.

The Laplace transform is a linear operation, since it's based on integration. Mathematically, 
\begin{align*}
	&x_1(t) \Longleftrightarrow X_1(s) &x_2(t) \Longleftrightarrow X_2(s)
\end{align*}

Example 4.1: For a signal \(x(t) = e^{-at}u(t)\), find the Laplace transform \(X(s)\) and its region of convergence (ROC).
\begin{align*}
	&X(s) = \int_{-\infty}^{\infty} e^{-at} u(t) e^{-st}dt \\
	&X(s) = \int_{0}^{\infty} e^{-at} e^{-st}dt = \int_{0}^{\infty}e^{-(s+a)t}dt \\
\end{align*}

\textbf{Important:} If the real value of \(s+a\) is positive, the integral will converge; but if \textit{not}, it won't converge. We'll need to note this in our solution.

\begin{align*}
	&e^{-at}u(t) \Longleftrightarrow \frac{1}{s+a} &\mathrm{ROC:} Re(s) > -a 
\end{align*}

Consider, for example, \(x(t) = -e^{-at}u(-t)\). In this case, the unit step is flipped along the \(y\)-axis (anti-causal). We can still find the Laplace transform for this signal, since the upper bound can be changed to \(0\), with a lower bound of \(-\infty\). However, we will get the \textit{same} Laplace transform for this signal as the prior one. That's why ROC is important, because it will distinguish these two signals when they are transformed back into the time domain, where they are different signals.

\textbf{NOTE:} fill in the above with more details.

In general, ROC is vital to any bilateral Laplace transforms for the above reasons. In the table of transform pairs that we'll be provided, ROC will be included. 

The unilateral Laplace transform of the second, anti-causal signal would simply be \(0\), so bilateral transforms are vital for anti-causal signals to yield any kind of useful result.

However, moest practical signals and systems are causal, and we may generalize the Laplace transforms for causal signals as 
\[
	X(s) = \int_{0^-}^{-\infty}x(t)e^{-st}dt,
\]
and no ROC specification is necessary.

This shouldn't come up too much since we'll be working largely with causal signals. Table 4.1 provides select \textit{unilateral} Laplace transform pairs, useful for causal signals.

Drill 4.1: By direct integration, find the Laplace transform \(X(s)\) and the region of convergence of \(X(s)\) for the gate function shown in Fig. 4.2(a). (Add the fig later, but it's a rectangular unit pulse from 0 to 2). 

We won't often use direct integration, but its worth reviewing. Working through the entire integration, 
\begin{align*}
	X(s) &= \int_{-\infty}^{\infty} [u(t) - u(t-2)]e^{-st}dt \\
			 &=\int_{-\infty}^{\infty} u(t) e^{st} dt - \int_{-\infty}^{\infty}u(t-2)e^{-st}dt \\
			 &=\int_{0}^{\infty}e^{st}dt - \int_{2}^{\infty}e^{-st}dt \\
			 &=\frac{1}{s}e^{-st} \Big|_0^\infty - [\frac{-1}{s}e^{-st}\Big|_2^\infty] \\
			 &=(0-\frac{-1}{s}e^0)-[0-\frac{-1}{s}e^{-2s}] \\
			 &=\frac{1}{s} - \frac{1}{s}e^{-2s}=\frac{1}{s}[1-e^{-2s}].
\end{align*}

In this case, direct integration was relatively simple (but still took multiple steps), and no ROC was needed since \(x(t)\) was a causal signal, but in most cases we'll use a table of transforms or software to carry out a Laplace transform.

Inverse Laplace transforms were covered in Circuits 2, but worth reviewing. If a direct transform pair can't be found in the table, partial fraction expansion of \(X(s)\) is generally used to break the s-valued function into partial fractions that are listed in the table of inverse transforms. (Review Example 4.3 in the textbook for a worked-through unilateral inverse Laplace transform problem using partial fraction decomposition. Probably worth reviewing PFD using the cover-up method, especially for complex-valued fractions).

Given a function \(X(s) = \dfrac{N(s)}{D(s)}\), the function is considered \textbf{strictly proper} if the degree of \(N(s)\) is less than the degree of \(D(s)\); if the degree of \(N(s)\) is equal to the degree of \(D(s)\), the function is considered \textbf{proper}. In the second case, we'll need an \(A\) constant in our partial fraction expansion (see section B.5 in the textbook for details. See Example 4.4 as well for examples of MATLAB-based partial fraction decomposition with strictly proper, proper, and improper functions \(X(s)\)).

MATLAB gives you three vectors: r, p, and k; the residues, poles, and polynomial coefficients, respectively. Residues correspond to the numerators of the partial fraction, poles to the zeroes of the denominator, and polynomial coefficients to the added constant(s). When running the residue() command in MATLAB, pay close attention to the order of the coefficients, and insert zeroes where there is no \(s\) value in the polynomial.

As an example, determine the inverse Laplace transform of 
\[
  X_b(s) = \frac{2s^2 + 7s +4}{(s+1)(s+2)^2}.
\]
Using MATLAB's conv() (\textbf{why?}) and residue() functions, we find 
\[
  X_b(s) = \frac{3}{(s+2)} + \frac{2}{(s+2)^2} - \frac{1}{(s+1)}.
\]

Using an example with complex roots, find the inverse Laplace transform of
\[
  \frac{8s^2 + 21s + 19}{(s+2)(s^2 + s + 7)}.
\]
MATLAB can handle complex values as well as real ones, and though we will find that both residues and poles are complex in this case, we may use angle() and abs() to express the complex values in polar form, which will then correspond to the table of inverse Laplace transforms.

Laplace and inverse Laplace transforms may also be found using MATLAB's symbolic math toolbox, a useful approach when dealing with complex functions that have no corresponding formula in a table of transforms. In MATLAB, this is done using laplace() and ilaplace().

Another example: Show that the Laplace transform of \(10e^{-3t} \cos{(4t + 53.13 \degree)}\) is \(\dfrac{(6s-14)}{(s^2+6s+25)}\). Use Table 4.1.

Using Table 4.1, 
\[
	re^{-at}\cos{(bt + \theta) \longleftrightarrow \frac{As+B}{s^2+2as+c}}.
\]

Starting with \(r=10\), from the table, \(r = \sqrt{\frac{A^2C+b^2-2ABa}{c-a^2}}\). Glancing at the right-hand side, \(A=6\), \(B=-14\), \(a=3\), \(c=25\). Therefore, 
\[
  r = \sqrt{\frac{6^2(25) - 2(6)(-14)(3)}{25-3^2}} = \sqrt{100} = 10.
\]
Then,
\[
  b=\sqrt{c-a^2} = \sqrt{25-3^2} = \sqrt{16} = 4.
\]
Again, using the table, 
\[
	\theta = \tan^-1{(\frac{6(3) + 14}{6 \ \sqrt{25-32}})} = 53.13 \degree.
\]

Extra practice using MATLAB residue() function.

\setcounter{subsection}{2}
\subsection{Solution of Differential and Integro-Differential Equations}
Beginning to look at \textit{applications} of the Laplace transform, one of its most important uses is in solving differential equations that would otherwise be time-consuming and/or difficult to solve in the time domain.

Differential and integro-differential equations arise naturally when capacitors and inductors are part of a circuit, since the behavior of both is linked to the time-derivatives of voltage and current, respectively; the Laplace transform converts this sort of behavior into a simple algebraic equation. For example, 
\[
	i_c(t) = C \frac{dV_c(t)}{dt} \to I_c(s) = C \cdot [sV_c(s) - v_c(0^-)].
\]

\textbf{Example 4.12}:

Solve 
\[
	(D^2 + 5D +6)y(t) = (d+1)x(t)
\]
for initial conditions \(y(0^-) = 2\) and \(\dot{y}(0^-) = 1\) and the input \(x(t) = e^{-4t}u(t)\).
Utilizing the time differentiation property of the Laplace transform,
\[
  \frac{dx(t)}{dt} = sX(s) - x(0^-),
\]
and converting \(x(t)\) to the s-domain,
\[
	[s^2Y(s) - 2s -1]+5[sY(s) -2]+6Y(s) = \frac{s}{s+4}+\frac{1}{s+4}.
\]
The above can be easily solved for \(Y(s)\), which may then be converted back to the time-domain. Solving, 
\[
  Y(s) = \frac{2s^2+20s+45}{(s+2)(s+3)(s+4)}.
\]
Using partial fraction decomposition to prepare for the use of the reverse Laplace transform,
\[
  Y(s) = \frac{13/2}{s+2}-\frac{3}{s+3}-\frac{3/2}{s+4},
\]
and thus
\[
	y(t) = (\frac{13}{2}e^{-2t}-3e^{-3t}-\frac{3}{2}e^{-4t})u(t).
\]

Put simply, we can easily convert a differential equation representing a system to the s-domain, solve the result algebraically, and then convert back to the time domain to find the system's total response to some input \(x(t)\). However, we can still extract the ZSR and ZIR from the result.

\textbf{Drill 4.7}
Solve 
\[
  \frac{d^2y(t)}{dt^2}+4\frac{dy(t)}{dt} +3y(t) = 2\frac{dx(t)}{dt}+x(t)
\]
for the input \(x(t) = u(t)\). The initial conditions are \(y(0^-)=1\) and \(\dot{y}(0^-)=2\).

Again, we'll convert to the s-domain to solve the system algebraically:
\[
	[s^2Y(s)-5y(0^-)-\dot{y}(0^-)]+4[sY(s)-y(0^-)]+3Y(s)=2[sX(s)-x(0^-)].
\]
Substituting intial conditions,
\[
	Y(s)[s^2+4s+3]+[-s-2-4]=(2s+1)X(s)-x(0^-).
\]
Converting other pieces to s-domain,
\begin{align*}
	&x(t)=u(t) \to X(s) = \frac{1}{s} \\
	&x(0^-) = 0
\end{align*}

Now,
\[
	Y(s)[s^2+4s+3]=(2+\frac{1}{s}+(s+6),
\]
and 
\[
  Y(s) = \frac{2+\frac{1}{s}+s+6}{s^2+4s+3} = \frac{s^2+8s+1}{s(s^2+4s+3)} = \frac{s^2+8s+1}{s(s+1)(s+3)}.
\]
Using partial fraction expansion,
\[
  Y(s) = \frac{\frac{1}{3}}{s}+\frac{3}{s+1}+\frac{\frac{-7}{3}}{s+3},
\]
and thus
\[
	y(t) = [\frac{1}{3} + 3e^{-t} - \frac{7}{3}e^{-3t}]u(t) = \frac{1}{3}(1+9e^{-t}-7e^{-3t})u(t).
\]

Similarly, the Laplace transform can be used to easily find the transfer function of some system, where 
\[
  H(s) \triangleq \frac{Y(s)}{X(s)}=\frac{P(s)}{Q(s)},
\]
with zero initial conditions. Brief reminder that \(H(s)\) is the Laplace transform of the impulse response, \(\mathcal{L}\{h(t)\}\). Similar to the impulse response, the transfer function allows us to find the system response to any input in the s-domain.

Unsurprisingly, the transfer function's roots are connected to system stability. Assuming \(H(s)\) is a proper function (\(M \leq N)\), the system is BIBO stable if and only if all the poles of \(H(s)\) are in the LHP. Otherwise, \((M > N)\), and the system is BIBO unstable as a rule.

\textbf{Note:} A pole is any value that causes the denominator to become zero, thereby bringing the function's value to infinity; in other words, a pole is any root of the polynomial found in the denominator of the transfer function.

In short, when designing systems, avoid improper transfer function (more poles than zeros, ideally, and at least as many poles as zeros).

A brief review of the approaches to stability analysis we've covered thus far. 

As an example, consider a simple system described by
\[
	y(t) = 2\dot{x}(t).
\]
Before doing any actual analysis, we could imagine that the input to the system is \(x(t) = u(t)\), in which case \(\dot{x}(t) = \delta(t)\), showing that the system is BIBO-unstable (the input is bounded, but the output is unbounded).

We may also take the analytical approach, and convert the system to the s-domain, as
\[
  Y(s) = 2sX(s).
\]
Finding the transfer function,
\[
  H(s) = \frac{Y(s)}{X(s)} = 2s = \frac{2s}{1}.
\]
Since the transfer function is \textit{improper}, the system is considered to be BIBO-unstable.

Another method of determining the system stability is to analyze the impulse response, found by taking the inverse Laplace transform of the transfer function, as
\[
	h(t) = \mathcal{L}^{-1}\{H(s)\} = \mathcal{L}^{-1}\{2s\} = 2\dot{\delta}(t).
\]
Then, we may determine whether the impulse response is absolutely integrable,
\[
	\int_{-\infty}^{\infty} |h(\tau)| d \tau = \int_{-\infty}^{\infty}|2\dot{\delta}(t)| = 2\delta(t).
\]

Since, by definition, the delta function \(\delta(t)\) is valued at infinity, the system does \textit{not} meet the criteria for BIBO-stability.

Stability analysis naturally applies to circuit analysis as well.

As an example, consider a simple series RC circuit. Can we use stability analysis to learn more about this circuit? 

Converting the circuit quantities to the s-domain, we can formulate the system's transfer function as
\[
  H(s) = \frac{I_o(s)}{V_i(s)} = \frac{Cs}{RCs+1}.
\]
Since the function is proper, we may assume the system is BIBO-stable.

Extending the analysis, consider the same circuit, this time with no resistance. Naturally, this system would be BIBO-\textit{unstable}, since an instantaneous jump in voltage would require an infinite amount of current (in other words, the time derivative of voltage would be \(\delta(t)\)).

In the s-domain, we can make similar conclusions using the transfer function,
\[
  H(s) = \frac{I_o(s)}{V_i(s)} = \frac{Cs}{1},
\]
and the improper function leads us to the same conclusion: the system is BIBO-unstable.

Note that by changing the definition of the system by instead considering the output to be \(v_i(t)\) rather than \(i_o(t)\), the "new" system becomes stable. The system's definition is critical to analyzing it; a circuit in and of itself is not a system, and its stability depends on the way in which we define the input/output relationship, and what parameters we choose to analyze. 

Side note, given a system \(S\), that system's inverse , \(S_i\), is defined by inverse of its transfer function, 
\[
  H_i(s) = \frac{1}{H(s)}.
\]

\subsection{Analysis of Electrical Networks: The Transformed Network}

As seen in the previous sections, one of the major benefits of the Laplace transform is converting complicated integro-differential equations to a simpler, algebraic form. 

Coincidentally, the mathematical description of most electrical circuits are defined as integro-differential equations; thus, learning to convert the description of those circuits to the s-domain will be useful.

Example 4.17: 

Find the loop current \(i(t)\) in the circuit shown in Fig. 4.10a if all the initial conditions are zero.

Converting the circuit's given parameters to the s-domain, 
\begin{align*}
	&10u(t) = \frac{10}{s} \\
	&Z(s) = s + 3 + \frac{2}{s} = \frac{s^2 + 3s + 2}{s},
\end{align*}
and so
\[
  I(s) = \frac{V(s)}{Z(s)} = \frac{10}{s^2+3s+2}.
\]

Then, working our way \textit{back} to the time domain, 
\[
	\mathcal{L}^{-1}\{I(s)\} = \mathrm{\textbf{fill this in}}.
\]
  
What about non-zero initial conditions? Considering an already charged capacitor in the above circuit, we may find the current at \(t=0^-\) by another Laplace transform,
\[
	I(s) = \mathcal{L}\{C \frac{dv(t)}{dt}\} = C[sV(s) - v(0^-)].
\]

Then, we can proceed, since the given initial voltage may be used to solve for \(I(s)\).

Example 4.13:

In the circuit of Fig. 4.7a, the switch it in the closed position for a long time before \(t=0\), when it is opened intantaneously. Find the inductor current \(y(t)\) for \(t \geq 0\).

First, we want to find initial conditions at \(t=0^-\). The switch is closed at this point, and we will have a voltage \(v_c(0^-)\) across the capacitor and a current \(y(0^-)\) in the loop, which we can easily derive as \(v_c(0^-) = 10 \ \mathrm{V}\) and \(y(0^-) = 2 \ \mathrm{A}\)m.

Now, we can transform the system to the s-domain as
\begin{align*}
	&10u(t) = \frac{10}{s} \\
	&V_c(s) = \frac{5}{s}, \ \frac{10}{s} \ \mathrm{\textbf{Fill this in}}.
\end{align*}

Then, using KVL, we may solve the circuit in the s-domain, and then convert back to time-domain, as
\begin{align*}
	Y(s) &= \frac{V}{Z} = \frac{10/s + 2 - 10/s}{s + 2 + 5/s} \\
			 &=\frac{2}{s+2+5/s} = \frac{2s}{s^2+2s+5}.
\end{align*}

The result of evaulating the denominator produces complex roots, and so we use 
\[
	re^{-at}cos(bt+\theta)u(t) = \frac{As+B}{s^2+2as +c}
\]
to find the time domain solution,
\[
	y(t) = \sqrt{5}e^{-t}\cos{(2t+26.6\degree)}u(t).
\]

\subsection{Block Diagrams}

Complex systems are best represented using the block diagram format, which simplifies the system to blocks (representing transfer functions) and inputs/outputs, as well as encoding the relationship between the many interconnected systems that usually make up a larger system. We may also include operations in the block diagram, representing addition, multiplication, addition, and differentiation.

Cascaded systems are constructed by connecting the output of one system to another, and reduce mathematically to \(H_1(s)H_2(s)\). Parallel connected systems sum the output of two systems to create one ouput, and reduce to \(H_1(s) + H_2(s)\).

Feedback systems are slightly more complicated, but very common. In feedback systems, the output of one system \(G(s)\) is processed by a second system, \(H(s)\), and \textit{added} to the input of \(G(s)\), generating a feedback loop. This sort of system interconnection reduces to 
\[
  \frac{G(s)}{1+G(s)H(s)}.
\]

Can we derive the transfer function of a negative feedback loop? (see the example in the slides/textbooks)

Recall the ouptut of any system is given by the convolution of its input and impulse response, 
\[
  y(t) = x(t) * h(t),
\]
or, in s-domain, and using the variables of our block diagram,
\[
  Y(s) = G(s)E(s).
\]
Since \(E(s)\) represents the sum of our feedback and the input signal \(X(s)\) (in actuality, this is a subtraction, but the sum symbol is used alongside a minus sign), we may expand this to
\begin{align*}
	Y(s) &= G(s)[X(s)-H(s)Y(s)] \\
			 &=G(s)X(s) - G(s)H(s)Y(s) \\
	Y(s)[1+G(s)H(s)] &= G(s)X(s) \\
	Y(s) &= \frac{G(s)}{1+G(s)H(s)}X(s).
\end{align*}

MATLAB example (see slides/textbook).

\subsection{System Realization}

Given a transfer function, we may develop a method to \textit{realize} the system described. In its most general form, a transfer function may be represented as
\[
	H(s) = \frac{b_0s^N + b_1s^{N-1} + ... + b_{N-1}s+b_N}{s^N + a_1s^{N-1} + ... + a_{N-1}s+a_N} = \frac{Y(s)}{X(s)}.
\]
Though there is no unique way of realizing a system, we may use integrators, adders, and multipliers to realize any arbitrarily chosen system with the form given above. Recall that in the s-domain,
\begin{align*}
	&\frac{d}{dt} \Leftrightarrow s \\
  &\int \Leftrightarrow \frac{1}{s}.
\end{align*}
Given an arbitrary transfer function, we could choose to split it into expressions that fit the above categories.
Example 4.22:
	Realize the transfer function \(H(s) = \frac{5}{s+7}\). Converting to the "canonical",  
\[
  H(s) = \frac{0s' + b_1}{s' + a_1},
\]
with 
\begin{align*}
 &b_0 = 0 \\
 &b_1=5 \\
 &a_1=7.
\end{align*}

Then, with this in mind, we may proceed to the block diagram realization, simplifying as needed. 

Partial fraction expansion, as usual, never fails; very useful for realizing parallel systems. (Sometimes we \textit{want} to implement a system in cascade or parallel form for one reason or another).

We'll talk about frequency response next time.

\setcounter{subsection}{7}
\subsection{Frequency Response of an LTIC System}

In the realm of analysis, the frequency response of a system (in most of our applications, a circuit) is a critical property of systems. As one might expect, converting a system to the frequency domain will make this sort of analysis much easier. Given a transfer function \(H(s)\), we may find the amplitude and phase response of a system by plotting \(H(s)\) as a function of \(j\omega\), as \(H(j\omega\) (essentially, limiting the independent variable of \(H(s)\) to the imaginary part of \(s\)). Both amplitude and phase as a function of frequency may be derived this way.

Example 4.28: 

Find and sketch the frequency response for \textbf{(a)} an ideal delay of \(T\) seconds, \textbf{(b)} an ideal differentiator, and \textbf{(c)} an ideal integrator.

\begin{enumerate}[label=\alph*)]
	\item By time shifting some input and converting to the frequency domain,
		\[
		  x(t) = 
		\]
	\item An ideal differentiator has a transfer function \(H(s) = s\), in the s-domain. Representing its frequency response, \(|H(j\omega)| = \omega\), representing a high-pass filter, and \(\angle{H(j\omega)} = \frac{\pi}{2}\), remaining constant.
	\item For an ideal integrator, the process is similar, and yields \(|H(j\omega)| = \frac{1}{\omega}\), representing a low-pass filter, with a constant phase angle of \(\angle{H(j\omega)} = -\frac{\pi}{2}\).
\end{enumerate}
\setcounter{subsection}{8}
\subsection{Bode Plots}
\end{document}
